% =====================================================================
\chapter{Polinomio de Taylor}
\label{tema:6}
% =====================================================================

La recta tangente es la mejor aproximación de una función por un polinomio de grado~1 cerca de un punto. El polinomio de Taylor generaliza esa idea a cualquier grado: es el polinomio que tiene en el punto las mismas derivadas que la función hasta un orden dado. El teorema de Taylor con resto de Lagrange dice cuánto se equivoca esa aproximación, lo que permite calcular valores de funciones con un error garantizado, calcular límites y clasificar puntos críticos. Se siguen el capítulo~20 de Spivak~\cite{spivak}, el capítulo~7 de Apostol~\cite{apostol}, la §5.3 de los apuntes de \emph{Analysis~I}~\cite{cambA1}, la §6.3 del segundo volumen de OpenStax~\cite{osc2} y las sesiones~23--26 del curso 18.01SC del MIT~\cite{mit1801}.

\section{El polinomio de Taylor}

\begin{definicion}[Polinomio de Taylor]\label{def:taylor}
Sea $f$ derivable $n$ veces en $a$. El \emph{polinomio de Taylor} de orden $n$ de $f$ en $a$ es
\[
P_{n,a}f(x)=\sum_{k=0}^{n}\frac{f^{(k)}(a)}{k!}(x-a)^{k}=f(a)+f'(a)(x-a)+\frac{f''(a)}{2}(x-a)^{2}+\dots+\frac{f^{(n)}(a)}{n!}(x-a)^{n},
\]
y el \emph{resto} es $R_{n,a}f(x)=f(x)-P_{n,a}f(x)$. Si $a=0$ se habla de polinomio de \emph{Maclaurin}.
\end{definicion}

\begin{proposicion}[Caracterización]\label{prop:taylor-unico}
$P=P_{n,a}f$ es el único polinomio de grado a lo sumo $n$ que cumple $P^{(k)}(a)=f^{(k)}(a)$ para $k=0,1,\dots,n$.
\end{proposicion}

\begin{proof}
Todo polinomio de grado $\le n$ se escribe de forma única como $P(x)=\sum_{k=0}^{n}c_k(x-a)^{k}$. Derivando $j$ veces, el término $c_k(x-a)^{k}$ da $\frac{k!}{(k-j)!}c_k(x-a)^{k-j}$ si $k\ge j$ y $0$ si $k<j$; en $x=a$ solo sobrevive $k=j$, luego $P^{(j)}(a)=j!\,c_j$. Así, $P^{(j)}(a)=f^{(j)}(a)$ si y solo si $c_j=f^{(j)}(a)/j!$.
\end{proof}

\begin{ejemplo}[Polinomios de Maclaurin básicos]\label{ej:maclaurin}
(a) $f(x)=e^{x}$: todas las derivadas valen $e^{0}=1$ en~0, luego $P_{n,0}(x)=\sum_{k=0}^{n}\frac{x^{k}}{k!}$.
(b) $\sin$: las derivadas en~0 son, cíclicamente, $0,1,0,-1$, luego $P_{2m+1,0}\sin(x)=x-\frac{x^{3}}{3!}+\frac{x^{5}}{5!}-\dots+(-1)^{m}\frac{x^{2m+1}}{(2m+1)!}$. Análogamente, $P_{2m,0}\cos(x)=1-\frac{x^{2}}{2!}+\dots+(-1)^{m}\frac{x^{2m}}{(2m)!}$.
(c) $\ln(1+x)$: $f^{(k)}(x)=(-1)^{k-1}(k-1)!/(1+x)^{k}$ para $k\ge1$, luego $f^{(k)}(0)/k!=(-1)^{k-1}/k$ y $P_{n,0}(x)=x-\frac{x^{2}}{2}+\frac{x^{3}}{3}-\dots+(-1)^{n-1}\frac{x^{n}}{n}$.
(d) $(1+x)^{\alpha}$, $\alpha\in\R$: $f^{(k)}(0)=\alpha(\alpha-1)\cdots(\alpha-k+1)$, luego $P_{n,0}(x)=\sum_{k=0}^{n}\binom{\alpha}{k}x^{k}$, con $\binom{\alpha}{k}=\frac{\alpha(\alpha-1)\cdots(\alpha-k+1)}{k!}$.
\end{ejemplo}

\section{El teorema de Taylor}

\begin{teorema}[Taylor con resto de Lagrange]\label{teo:taylor}
Sea $f$ derivable $n+1$ veces en un intervalo abierto $I$ que contiene a $a$. Para cada $x\in I$ existe $\xi$ entre $a$ y $x$ tal que
\[
f(x)=\sum_{k=0}^{n}\frac{f^{(k)}(a)}{k!}(x-a)^{k}+\frac{f^{(n+1)}(\xi)}{(n+1)!}(x-a)^{n+1}.
\]
\end{teorema}

\begin{proof}
Si $x=a$ no hay nada que probar. Sea $x\ne a$ fijo y, para $t$ en el intervalo cerrado de extremos $a$ y $x$, sea
\[
g(t)=f(x)-\sum_{k=0}^{n}\frac{f^{(k)}(t)}{k!}(x-t)^{k}-K\frac{(x-t)^{n+1}}{(n+1)!},
\]
donde la constante $K$ se elige para que $g(a)=0$ (es posible porque $(x-a)^{n+1}\ne0$). También $g(x)=f(x)-f(x)=0$, pues en $t=x$ solo sobrevive el término $k=0$ del sumatorio. Al derivar respecto de $t$, cada término del sumatorio da, por la regla del producto, $\frac{f^{(k+1)}(t)}{k!}(x-t)^{k}-\frac{f^{(k)}(t)}{(k-1)!}(x-t)^{k-1}$ (sin el segundo sumando si $k=0$); la suma es telescópica y queda
\[
g'(t)=-\frac{f^{(n+1)}(t)}{n!}(x-t)^{n}+K\frac{(x-t)^{n}}{n!}.
\]
Por el teorema de Rolle aplicado a $g$ entre $a$ y $x$, existe $\xi$ estrictamente entre ellos con $g'(\xi)=0$; como $x-\xi\ne0$, se obtiene $K=f^{(n+1)}(\xi)$. Sustituyendo en $g(a)=0$ se llega a la fórmula del enunciado.
\end{proof}

Para $n=0$ el teorema es exactamente el teorema del valor medio: $f(x)=f(a)+f'(\xi)(x-a)$. El punto $\xi$ es desconocido, pero basta conocer una cota de $f^{(n+1)}$ para acotar el error.

\begin{corolario}[Cota del error y forma de Peano]\label{cor:cota}
Si además $\abs{f^{(n+1)}(t)}\le M$ para todo $t$ entre $a$ y $x$, entonces
\[
\abs{R_{n,a}f(x)}\le\frac{M}{(n+1)!}\abs{x-a}^{n+1}.
\]
En particular, si $f^{(n+1)}$ está acotada cerca de $a$, entonces $R_{n,a}f(x)/(x-a)^{n}\to0$ cuando $x\to a$, lo que se escribe $f(x)=P_{n,a}f(x)+o\bigl((x-a)^{n}\bigr)$.
\end{corolario}

\begin{proof}
La cota es inmediata a partir del \cref{teo:taylor}. Para la segunda parte, $\abs{R_{n,a}f(x)/(x-a)^{n}}\le\frac{M}{(n+1)!}\abs{x-a}\to0$.
\end{proof}

La notación $o(g(x))$ («o pequeña») designa cualquier función $h$ con $h(x)/g(x)\to0$; se cumplen $o(x^{n})+o(x^{n})=o(x^{n})$, $x^{m}o(x^{n})=o(x^{n+m})$ y $o(x^{m})=o(x^{n})$ si $m\ge n$, cuando $x\to0$.

\begin{table}[h]
\centering\small
\caption{Desarrollos de Maclaurin más usados, cuando $x\to0$.}\label{tab:maclaurin}
\begin{tabular}{@{}ll@{}}
\toprule
$e^{x}=1+x+\dfrac{x^{2}}{2}+\dfrac{x^{3}}{6}+o(x^{3})$ & $\ln(1+x)=x-\dfrac{x^{2}}{2}+\dfrac{x^{3}}{3}+o(x^{3})$ \\[2mm]
$\sin x=x-\dfrac{x^{3}}{6}+o(x^{4})$ & $\cos x=1-\dfrac{x^{2}}{2}+\dfrac{x^{4}}{24}+o(x^{5})$ \\[2mm]
$(1+x)^{\alpha}=1+\alpha x+\dfrac{\alpha(\alpha-1)}{2}x^{2}+o(x^{2})$ & $\dfrac{1}{1-x}=1+x+x^{2}+x^{3}+o(x^{3})$ \\[2mm]
$\tanh x=x-\dfrac{x^{3}}{3}+o(x^{4})$ & $\sigma(x)=\dfrac12+\dfrac{x}{4}-\dfrac{x^{3}}{48}+o(x^{4})$ \\
\bottomrule
\end{tabular}
\end{table}

En la tabla, $\sin x=x-x^{3}/6+o(x^{4})$ (y no solo $o(x^{3})$) porque el término de grado~4 es nulo: $P_{3,0}\sin=P_{4,0}\sin$. Lo mismo ocurre con $\cos$, $\tanh$ y $\sigma-\tfrac12$, que son pares o impares.

\begin{ejemplo}[Desarrollos de $\tanh$ y de $\sigma$]\label{ej:taylor-activaciones}
Con $t=\tanh x$, $t'=1-t^{2}$, $t''=-2t(1-t^{2})$ y $t'''=-2(1-t^{2})(1-3t^{2})$ (\cref{eje:4-sucesivas}). En $x=0$, $t=0$, luego $t'(0)=1$, $t''(0)=0$ y $t'''(0)=-2$: $P_{3,0}\tanh(x)=x-\frac{2}{6}x^{3}=x-\frac{x^{3}}{3}$. Como $\sigma(x)=\tfrac12+\tfrac12\tanh(x/2)$ (\cref{eje:1-sigmoide-simetria}), $\sigma(x)=\tfrac12+\tfrac12\bigl(\tfrac x2-\tfrac{x^{3}}{24}\bigr)+o(x^{4})=\tfrac12+\tfrac x4-\tfrac{x^{3}}{48}+o(x^{4})$. Cerca del origen, $\sigma$ y $\tanh$ son prácticamente lineales; la no linealidad aparece en el término cúbico.
\end{ejemplo}

\section{Aproximación con error controlado}

\begin{ejemplo}[Cálculo de $e$]\label{ej:calculo-e}
Con $f=\exp$, $a=0$, $x=1$: $e=\sum_{k=0}^{n}\frac{1}{k!}+\frac{e^{\xi}}{(n+1)!}$ con $0<\xi<1$. Como $e^{\xi}<e<3$, el error es menor que $3/(n+1)!$. Para que sea menor que $10^{-6}$ basta $(n+1)!>3\cdot10^{6}$; como $9!=362\,880$ y $10!=3\,628\,800$, basta $n=9$: $\sum_{k=0}^{9}1/k!=2{,}7182815\dots$, y $e=2{,}7182818\dots$
\end{ejemplo}

\begin{ejemplo}[Aproximación de una raíz]\label{ej:raiz}
Se aproxima $\sqrt{1{,}1}$ con $f(x)=\sqrt{1+x}$ en $a=0$ y $n=1$: $P_{1,0}f(x)=1+x/2$, y $f''(t)=-\tfrac14(1+t)^{-3/2}$. Para $x=0{,}1$, $\sqrt{1{,}1}\approx1{,}05$ con $\abs{R}\le\frac{1}{2}\cdot\frac14\cdot(0{,}1)^{2}=0{,}00125$, porque $(1+t)^{-3/2}\le1$ para $t\ge0$. Además el resto es negativo ($f''<0$), luego $1{,}05$ es una aproximación por exceso. El valor exacto es $1{,}048808\dots$, con error $0{,}00119$.
\end{ejemplo}

\section{Aplicaciones a límites y a extremos}

\begin{ejemplo}[Límites mediante desarrollos]\label{ej:limites-taylor}
(a) $\dfrac{x-\sin x}{x^{3}}=\dfrac{x-(x-x^{3}/6+o(x^{4}))}{x^{3}}=\dfrac16+o(x)\to\dfrac16$, como en el \cref{ej:lhopital}(a), pero sin derivar tres veces.
(b) $\dfrac{e^{x}-1-x}{x^{2}}=\dfrac{x^{2}/2+o(x^{2})}{x^{2}}\to\dfrac12$.
(c) $\dfrac{\ln(1+x)-x}{x^{2}}=\dfrac{-x^{2}/2+o(x^{2})}{x^{2}}\to-\dfrac12$.
\end{ejemplo}

\begin{proposicion}[Criterio de la derivada $n$-ésima]\label{prop:criterio-n}
Sea $f$ de clase $C^{n}$ en un entorno de $a$, con $f'(a)=\dots=f^{(n-1)}(a)=0$ y $f^{(n)}(a)\ne0$, $n\ge2$. Si $n$ es par, $f$ tiene en $a$ un mínimo local estricto si $f^{(n)}(a)>0$ y un máximo local estricto si $f^{(n)}(a)<0$. Si $n$ es impar, $f$ no tiene extremo local en $a$.
\end{proposicion}

\begin{proof}
Por el \cref{teo:taylor} de orden $n-1$, $f(x)-f(a)=\frac{f^{(n)}(\xi_x)}{n!}(x-a)^{n}$ con $\xi_x$ entre $a$ y $x$. Por la continuidad de $f^{(n)}$ y la \cref{prop:signo}, $f^{(n)}(\xi_x)$ tiene el signo de $f^{(n)}(a)$ si $x$ está suficientemente próximo a $a$. Si $n$ es par, $(x-a)^{n}>0$ para $x\ne a$ y $f(x)-f(a)$ tiene el signo de $f^{(n)}(a)$. Si $n$ es impar, $(x-a)^{n}$ cambia de signo en $a$ y $f(x)-f(a)$ también.
\end{proof}

\begin{observacion}[Errores frecuentes]\label{obs:6-errores}
(a) Evaluar las derivadas en $x$ en lugar de en $a$: los coeficientes del polinomio de Taylor son números, $f^{(k)}(a)/k!$.
(b) Olvidar el factorial, o escribir $(x-a)^{k}$ como $x^{k}-a^{k}$.
(c) Tratar $\xi$ como un número conocido o como si fuera igual a $a$; solo se sabe que está entre $a$ y $x$, y por eso se usa una cota de $f^{(n+1)}$ en todo ese intervalo.
(d) Usar el polinomio de Taylor lejos de $a$: $P_{n,0}$ de $\ln(1+x)$ no aproxima $\ln(1+x)$ para $x>1$, por grande que sea $n$, y en general el error crece como $\abs{x-a}^{n+1}$.
(e) Truncar los desarrollos a distinto orden en el numerador y el denominador de un límite: hay que llegar, en ambos, hasta el primer término no nulo.
\end{observacion}

\begin{observacion}[Conexión con el curso]
El desarrollo de Taylor de segundo orden, en varias variables, es la herramienta del módulo~4 para clasificar puntos críticos con el hessiano y para analizar en el módulo~7 el descenso de gradiente, que se basa en la aproximación $L(\theta+h)\approx L(\theta)+L'(\theta)h+\tfrac12L''(\theta)h^{2}$. El mismo teorema da el error de las diferencias finitas con que se comprueban los gradientes (módulo~4, \cref{eje:6-diferencias}). En el módulo~11, $\ln(1+x)\approx x$ para $x$ pequeño explica por qué las bibliotecas numéricas ofrecen una función específica para $\ln(1+x)$, y en el módulo~9 la aproximación de la GELU mediante $\tanh$ con un término cúbico es un ajuste del mismo tipo que los de este tema.
\end{observacion}

\begin{lecturas}
\item Spivak~\cite{spivak}, cap.~20 «Approximation by Polynomial Functions».
\item Apostol~\cite{apostol}, cap.~7 «Polynomial Approximations to Functions».
\item Gowers y Chua~\cite{cambA1}, §5.3 «Differentiation theorems» (teorema de Taylor con resto de Lagrange).
\item OpenStax~\cite{osc2}, §6.3 «Taylor and Maclaurin Series».
\item MIT 18.01SC~\cite{mit1801}, sesiones 23--26 (aproximaciones lineal y cuadrática).
\end{lecturas}

\section*{Ejercicios del tema 6}
\addcontentsline{toc}{section}{Ejercicios del tema 6}

\begin{ejercicio}{1}{eje:6-coeficientes}
(a) Se sabe que $P_{2,0}f(x)=3-x+5x^{2}$. Halle $f(0)$, $f'(0)$ y $f''(0)$. ¿Tiene $f$ un extremo local en~0? (b) ¿Cuál es el polinomio de Taylor de orden~3 en $a=1$ de $p(x)=x^{3}$? Escríbalo en potencias de $(x-1)$ y compruebe que coincide con $p$.
\end{ejercicio}

\begin{ejercicio}{1}{eje:6-por-que}
Explique por qué el polinomio de Taylor de orden~$n$ de un polinomio de grado $\le n$ coincide con él, por qué el de una función par en $a=0$ solo tiene potencias pares y por qué $P_{3,0}\sin=P_{4,0}\sin$.
\end{ejercicio}

\begin{ejercicio}{2}{eje:6-maclaurin}
Calcule a partir de la definición: (a) $P_{4,0}$ de $\cos x$; (b) $P_{3,0}$ de $e^{-x}$; (c) $P_{3,0}$ de $\dfrac{1}{1-x}$; (d) $P_{2,0}$ de $\sqrt{1+x}$; (e) $P_{3,0}$ de $\arctan x$.
\end{ejercicio}

\begin{ejercicio}{2}{eje:6-otros-puntos}
Calcule: (a) $P_{3,1}$ de $\ln x$; (b) $P_{2,4}$ de $\sqrt{x}$; (c) $P_{3,1}$ de $1/x$. Use (a) para aproximar $\ln 1{,}1$ y compare con el valor real.
\fuente{apartado (a) adaptado de OpenStax~\cite{osc2}, §6.3, ejemplo~6.11.}
\end{ejercicio}

\begin{ejercicio}{2}{eje:6-ln12}
(a) Escriba la aproximación cuadrática general de una función dos veces derivable cerca de $x=a$. (b) Úsela para aproximar $\ln 1{,}2$. (c) Acote el error con el resto de Lagrange y compruebe la cota con el valor real.
\fuente{adaptado de MIT 18.01SC, examen~2~\cite{mit1801ex}, problema~1.}
\end{ejercicio}

\begin{ejercicio}{2}{eje:6-cuadraticas}
Halle las aproximaciones cuadráticas cerca de $x=0$ de: (a) $\tan x$; (b) $\ln\cos x$; (c) $\sin 2x$ y $\cos 2x$; (d) $\dfrac{1}{\cos x}$; (e) $e^{x^{2}}$; (f) $\dfrac{1}{a+bx}$ con $a\ne0$.
\fuente{adaptado de MIT 18.01SC, Jerison y Mattuck~\cite{mit1801ex}, ejercicios 2A-6, 2A-12(d) y 2A-13.}
\end{ejercicio}

\begin{ejercicio}{2}{eje:6-error}
(a) Aproxime $\sin 0{,}1$ con $P_{3,0}\sin$ y dé una cota del error. (b) ¿Qué orden $n$ de $P_{n,0}$ de $e^{x}$ garantiza $e^{0{,}5}$ con error menor que $10^{-5}$? (c) Aproxime $\sqrt[3]{8{,}1}$ con un polinomio de Taylor de orden~1 de $\sqrt[3]{x}$ en $a=8$ y acote el error.
\end{ejercicio}

\begin{ejercicio}{2}{eje:6-limites}
Calcule con desarrollos de Taylor los límites cuando $x\to0$ de:
\begin{enumerate}[label=(\alph*)]
  \item $\dfrac{e^{x}-1-x-x^{2}/2}{x^{3}}$;
  \item $\dfrac{\cos x-1+x^{2}/2}{x^{4}}$;
  \item $\dfrac{\sin x-x\cos x}{x^{3}}$;
  \item $\dfrac{\tanh x-x}{x^{3}}$;
  \item $\dfrac{\sigma(x)-\tfrac12-\tfrac x4}{x^{3}}$.
\end{enumerate}
\end{ejercicio}

\begin{ejercicio}{2}{eje:6-poligono}
(a) Con el teorema del coseno, demuestre que el lado de un polígono regular de $n$ lados inscrito en la circunferencia unidad mide $\sqrt{2-2\cos(2\pi/n)}$. (b) Con la aproximación cuadrática del coseno, calcule el límite del perímetro cuando $n\to\infty$. (c) Calcule el perímetro para $n=6$ y $n=100$ y compare con $2\pi$.
\fuente{adaptado de MIT 18.01SC, Jerison y Mattuck~\cite{mit1801ex}, ejercicio 2A-16.}
\end{ejercicio}

\begin{ejercicio}{3}{eje:6-taylor-orden1}
Sea $f$ dos veces derivable en un intervalo abierto que contiene a $a$ y a $x$. Demuestre, aplicando dos veces el teorema del valor medio de Cauchy a $F(t)=f(t)-f(a)-f'(a)(t-a)$ y $G(t)=(t-a)^{2}$, que existe $\xi$ entre $a$ y $x$ con $f(x)=f(a)+f'(a)(x-a)+\tfrac12f''(\xi)(x-a)^{2}$. Deduzca que si $f''\ge0$, la gráfica está por encima de la tangente.
\end{ejercicio}

\begin{ejercicio}{3}{eje:6-binomica}
(a) Enuncie el teorema de Taylor. (b) Sea $c_n=\binom{1/2}{n}$. Demuestre que $\abs{c_n}\le1$ para todo $n\ge0$. (c) Demuestre que, para $x\in(0,1)$, $\sqrt{1+x}=\lim_{N\to\infty}\sum_{n=0}^{N}c_nx^{n}$, acotando el resto de Lagrange.
\fuente{adaptado de Cambridge, \emph{Tripos} IA 2013, Paper~1~\cite{tripos2013}, pregunta~9D(b).}
\end{ejercicio}

\begin{ejercicio}{3}{eje:6-ln2}
Demuestre, acotando el resto de Lagrange, que para $x\in[0,1]$ se cumple $\abs{\ln(1+x)-P_{n,0}(x)}\le\dfrac{x^{n+1}}{n+1}$. Deduzca que $\ln2=\lim_{n\to\infty}\Bigl(1-\dfrac12+\dfrac13-\dots+\dfrac{(-1)^{n-1}}{n}\Bigr)$ y diga cuántos términos garantizan un error menor que $10^{-3}$.
\end{ejercicio}

\begin{ejercicio}{3}{eje:6-e-irracional}
Demuestre que $e$ es irracional. (Indicación: si $e=p/q$ con $p,q\in\N$, multiplique la igualdad del \cref{ej:calculo-e} con $n=q$ por $q!$ y observe que el resto multiplicado por $q!$ tendría que ser un entero estrictamente entre $0$ y $1$.)
\end{ejercicio}

\begin{ejercicio}{3}{eje:6-criterio-n}
(a) Aplique la \cref{prop:criterio-n} para clasificar el punto crítico $x=0$ de: $x^{4}$, $x^{5}$, $\cos x-1+\dfrac{x^{2}}{2}$ y $x\sin x$. (b) Demuestre que $\sin x\ge x-\dfrac{x^{3}}{6}$ para todo $x\ge0$ y que $\cos x\ge1-\dfrac{x^{2}}{2}$ para todo $x\in\R$, usando el signo del resto de Lagrange.
\end{ejercicio}

\begin{ejercicio}{3}{eje:6-diferencias}
Sea $f$ de clase $C^{3}$ con $\abs{f''}\le M_2$ y $\abs{f'''}\le M_3$. Demuestre que
\[
\Bigl|\frac{f(a+h)-f(a)}{h}-f'(a)\Bigr|\le\frac{M_2}{2}\abs h,\qquad
\Bigl|\frac{f(a+h)-f(a-h)}{2h}-f'(a)\Bigr|\le\frac{M_3}{6}h^{2}.
\]
Compruebe numéricamente las dos cotas con $f=\sin$, $a=1$ y $h=0{,}01$, e indique cuál de las dos fórmulas es preferible para comprobar gradientes.
\end{ejercicio}
